\documentclass[kvant.tex]{subfiles}

\begin{document}
1. Вектор r\textsubscript{1}, проведенный из точки m\textsubscript{1} в центр вращения, есть сумма двух векторов, направленных от m\textsubscript{1} к m\textsubscript{2} и от m\textsubscript{1} к m\textsubscript{3}. Модули этих векторов равны \(\frac{m_2}{m_1+m_2+m_3}R\) и \(\frac{m_3}{m_1+m_2+m_3}R\) (см. рис. 1).\par
Параллелограмму с такими сторонами подобен параллелограмм сил \(F_{12} = \gamma\frac{m_1m_2}{r^2}\), ибо сумма этих сил F\textsubscript{1}должна быть направлена к центру вращения. Из рассмотрения подобия
\[F_1: r_1 = \gamma\frac{m_1m_2}{R^2}:\frac{m_2}{m_1+m_2+m_3}R =\]\[ = \gamma\frac{m_1m_3}{R^2}:\frac{m_3}{m_1+m_2+m_3}R.\]
Отсюда
\[\omega^2 = \frac{a_1}{r_1} = \frac{F_1}{m_1r_1} = \gamma\frac{m_1+m_2+m_3}{R^3}.\]\par
2. При заданном в условии соотношении скоростей и их направлениях за малый про-межуток времени треугольних масс перейдет в подобный начальному. Если и для нового треугольника соотношение скоростей удов летворяет условию, то все будет доказано. Но так как ускорение пропорционально расстоянию до центра вращения (см. упражне вне 1), то и приращение скоростей, и новые скорости пропорциональны этим расстоя ниям, в направления скоростей изменятся горостей изменятся на один и тот же угол (см. рис. 2).
\end{document}